\section{The Freudian Framework: Dreams as Wish-Fulfillment}\label{sec:freud}

The scientific study of dreams is a recent development. For the first half of the 20th century, the field was dominated by the psychoanalytic theories of Sigmund Freud. Freud proposed that dreams were not random but were the ``royal road to the unconscious'', representing ``unconscious wishes, urges, and feelings'' \citep{freud1900dreams}. His wish-fulfillment model held that dreams have a \textit{manifest content} (the literal story and images) which serves to disguise a deeper \textit{latent content} (the hidden, symbolic, and often unacceptable meaning). He argued that dreams function as ``guardians of sleep'' by allowing the mind to ``do away with the wish'' in a hallucinatory experience, thus preventing the drive from waking the sleeper \citep{freud1900dreams}.

Freud's method of dream interpretation involved analyzing the symbols and narratives in dreams to uncover their latent content. He believed that many dream symbols were universal, such as snakes representing phallic symbols or houses symbolizing the self \citep{freud1900dreams}. However, Freud also acknowledged that individual experiences and associations could influence dream symbolism.

Carl Jung, a contemporary, broke with Freud but remained within the psychoanalytic framework, viewing dreams as direct mental expressions, compensations for waking life, and expressions of a ``collective unconscious'' through mythic archetypes \citep{jung1974dreams}. Until the 1950s, dream research remained entirely subjective and interpretive, operating within this psychoanalytic paradigm.

Ppsychoanalytic theories were groundbreaking and laid the foundation for the psychoanalytic approach to dream analysis. However, they have also been widely criticized for their lack of empirical support and scientific rigor. Modern sleep research has largely moved away from Freud's psychoanalytic framework, focusing instead on the neurobiological and cognitive aspects of dreaming. Nevertheless, Freud's work remains historically significant and continues to influence certain therapeutic practices and cultural understandings of dreams.

